\documentclass[12pt]{article}

\usepackage[a4paper, total={6in, 8in}]{geometry}
\usepackage{textcomp}

\begin{document}
\setlength{\parindent}{0pt}

\def \t {\textrightarrow}
\def \v {\vspace{1em}}
\def \bi {\begin{itemize}}
\def \ei {\end{itemize}}
\def \s[#1] {\section*{#1}}
\def \ss[#1] {\subsection*{#1}}
\def \sss[#1] {\subsubsection*{#1}}

\s[Governi di centro in italia]
Governi che permettono di superare il monocolore del governo italiano \\
Governi che dovrebbero essere maggiormente aperti di quelli monocolore di De Gasperi \t che pure aveva catturato molti voti dei conservatori, dichiara di non voler fare governi troppo conservatori \\
Con questa idea vengono varate delle riforme, che pero non tengono il passo

\ss[Rifoerma agraria]
Sono necessarie, perche al sud ci sono manifestazioni e occupazioni di braccianti al sud \t erano gia iniziate nel 48 \\
Nel 49 pero i contadini del sud intensificano proteste, perche vogliono distruggere il latifondo \t polizia dovra intervenire, che fermera l'occupazione delle terre \t due braccianti saranno uccisi in calabria \\
Riforma è necessaria \t e dice che bisogna espropriare i grandi latifondi del sud per ridistribuire terre ai contadini \\
Governo vuole cosi generare tante piccole imprese agricole, per dare vigore alla produzione diversificata e per dare + forza alla classe contadina, che diventava proprietaria anche \\
"Col diretti" è un organizzazione che doveva coordinare, ed era legata alla demcorazia cristiana \\
Questo avviene solo in parte \t le nuove aziende, per quanto piccole, saranno poco vive \t scara produttivita \\
Inoltre in sicilia vengono censite 67 mila famiglie idonee a ricevere appezzamento di terra, di cui sono 17 mila la riceveranno

\ss[Cassa del mezzogiorno]
Nel 50 viene istituita la cassa per il mezzogiorno \t ente amministrativo che doveva finanziare la rinascita del meridione, attraverso grandi opere pubbliche \\
In 10 anni devono essere realizzati ferrrovie, acquedotti etc. \t infrasttruture + agricoltura moderna dovevano garantire il decollo dell'industria nel meridione \\
La cassa pero degenera, e diventa uno strumento di appoggio del clientelismo del sud italia \t il risultato è fallimentare, come per la rif. agraria

\ss[Nascita della comunita europea]
È voluta da De Gasperi \t lui e Schumann e Adenauer sono considerati i padri dell ue

\v

Italia cresce, ma non in modo omogeneo e non risolve tutti i problemi \\
Il disagio si esprimera in scioperi e manifestazioni sostenuti dalla CGL e partiti di sinistra \t e le manifestazioni verranno represse da Mario Scelba, che sara ministro degli interni tra ? e ? \\
Con lui si crea la celere = reparto di polizia motorizzato, che si scontra spesso \t caso + grave a Modena, quando vengono uccisi 6 operai dalla polizia \\
Anche all'interno della maggioranza si crea dissenso \t sia sulle riforme, sia sul mantemnimeto dell'ordine sociale \t nel ? il partito liberale esce dal governo, che quindi perde forza di coalizione e si indebolisce \\
Nelle amministrative tra il 51 e ?? \t d.c. perde molti voti, mentre destre e sinistre crescono

\ss[La legge truffa]
De Gasperi pero vuole governo stabile e autorevole \t cosi nel 53 viene approvata una nuova legge elettorale, che il parlamento non voleva \t definita la "Legge truffa" \\
Lo schieramento di partiti che avesse ottenuto + 50 per cento dei voti avrebbe ottenuto il 65 per centro dei seggi in parlamento \t ricorda un po la legge Acerbo \\
Cosi si limita moltissimo l'opposizione \t il parlamento cosi non si divide + in modo equo \\
Questa legge doveva dare + stabilita al governo, come voleva De Gasperi \t doveva limitare l'opposizione, e rendere il governo anche + autoritario \\
Nel 53 non c'era possibilita di accordo tra maggioranza ed opposizione \t partiti di sinistra sapevano che non avrebbero mai ottenuto il 50 per cento \\
De Gasperi vuole semplicemente tutelare il suo partito

\v

Con le elezioni giugno 53, i partiti di governo non ottegono il 50 per cento \t la d.c. perde 2 milioni di voti, e a livello politico crescono le destre e sinistre (ora anche a livello amministrativo) \\
1953 segna la fine di De Gasperi come presidente del consiglio \t nuova pagina \t ci saranno ancora governi democristiani, ma senza De Gasperi e con quadro politico + instabile \\
Questo perche l italia sta cambiando \t all inizio anni 50 è ancora arretrato, ma a partire dal 48 ci sara boom economico \\
Questo governo non fa riforme sociali, ne di ampio respiro a livello internazionale \t non tengono passo con la crescita \\
Si sente l'esigenza di allargare la maggioranza parlamentare \t all'interno della d.c. si tendera + a destra, e si faranno entrare forze conservatrici (a volte anche autoritarie) \t es. nel 60 governo monocolore con ministri solo della d.c., ma con l'appoggio dell MSI \\
Si scatena protesta antifascita \t repressa, il governo Tambroni poi cade \t nella d.c. prendo quindi sopravvento forze + di sinistra \\
Il partito socialista non guarda + verso l'urss, non + filosovietico \t mentre comunisti si \\
D.c. vuole allargare maggioranza, e guarda verso partito socialista che è piu moderato 

\v

Questo accade in pieno boom economico \t crescita miracolosa del paese, italia molto industrializzata \t nel 58 italia entra nella CEE \\
Ragioni miracolo economico:
\bi
  \item ci sono molti imprenditori che sono creativi
  \item grande lavoro di operai e immigrati, che vanno da sud a nord
  \item grandi personaggi come Enrico Mattei
  \item inoltre ripresa dell'economia mondiale
  \item inoltre in italia c'è un eccedenza di offerta di mano d'opera \t forte disoccupazione, e questo abbassa il prezzo della mano d'opera \t imprenditori possono tenere salari bassi (tutti vogliono lavorare e sindacati poco presenti)
\ei
Settori + avanzati sono automobilistico, siderurgico e ?? \t italia compete nel mercato europeo

\v

Governi centristi quindi favoriscono infrasttruture necessarie per sviluppo \t ferrovie, e autostrade \t l'automobile è il simbolo del benessere, insieme agli elettrodomestici \\
Nel 55 c'è un automobile ogni 77 abitanti, mentre nel 57 ogni 39 \t tutti comprano macchina, facendo anche prestiti \\
Arriva anche il frigorifero \t prima c'era la ghiacciaia, che rivoluziona le abitudini \t arrivano anche i supermercati (nel 57 viene aperto a Roma il primo supermercato) \\
Arriva anche la lavatrice, prima c'erano le lavandaie \t arriva anche spirapolvere \\
Arriva anche il bagno nelle case \t prima erano nel pianerottolo, e il bagno era solo per case + di lusso \\
Arriva anche il cinema americano, che era vietato nel fascismo \t arrivano anche Rossellini, Fellini: i grandi registi \\
Arriva anche la televisione: la RAI inizia trasmissioni nel 54, che ha avuto un ruolo sociale molto importante \t tv ha insegnato la lingua italiana, c'erano anche dei programmi dedicati

\v

Questo pero non accade in tutta la penisola \t il nord diventa molto sviluppato, mentre il sud rimane arretrato \\
Quindi c'è una forte immigrazione interna \t in particolare verso i luoghi dove ci sono possibilita lavorative, in particolare Torino dove c'era la FIAT \\
Questo modifica anche l'assetto urbanistico di Torino \t nascono nuovi quartieri, e città si espande \\
C'è anche migrazione dalla campagna alla citta \t la lombardia va a milano, e anche in veneto si vede questo trend \\
Il triangolo industriale nasce ora \t qua si concentrano emigrati, e diventa zona + produttiva d'italia \\
A roma vanno abitanti dell'abruzzo e del molise \t e questo confluiva nel terziario, mentre a nord verso l'industria \\
Toscana rimane equilibrata \t riceve tanti emigrati quanti ne escono \t mentre la Sicilia perde un sacco di abitanti

\v

Conseguenza del boom economico sul piano politico? \t cambiamento va orientato, e squilibrio sociale va equilibrato \t se ne occupa il centro-sinistra \\
Partito socialista, dopo Ungheria, si avvicina a d.c. e si allonta da URSS \t e nella d.c. si afferma Aldo Moro, che sara l'interlocutore dei socialisti \t e sara il maggiore artefice di alleanze di centro sinistra \\
(1) Il papa Giovanni 23 ammorbidisce la posizione della chiesa nei confronti del socialismo, e questo favorisce il dialogo tra d.c. e partito socialista \\
(2) Contesto internazionale favorisce posizioni + riformiste del partito socialista e anche grazie ad (3) Aldo Moro \\
Inoltre (4) Kennedy è presidente degli USA nel 51 \t è favorevole al centro sinistra, perche pensava che potesse avviare riforme indispensabili per italia, e cosi facendo si emargina anche il partito comunista

\ss[Governo Fanfani]
Nel 52 entra la d.c. a Napoli \t Moro si afferma, che voleva collaborazione con socialisti \t nasce nuovo governo, guidato da Fanfani \t è il quarto governo fanfani, che è di colazione: d.c., partito repubblicano e social democratico \t riceve appoggio esterno dell partito socialista (che non vota contro) \\
Nasce quindi un governo di centro-sinistra, che si propone una serie di riforme (molte delle quali non verranno realizzati) \\
Fanfani fa alcune riforme:
\bi
  \item riforma della scuoal obbligo scolastico a 14 anni e ??
  \item nasce l'enel: nazionalizzata l'energia elettrica
\ei
Vero problema saranno gli enti + conservatori italiani \t es. gli enti amministrativi, e saranno l'ostacolo maggiore \t Fanfni voleva istituire le regioni, ma trovera difficolta \\
Il suo governo dura poco \t nel 63 sara aldo moro a formare il primo governo di cui fanno parte i socialisti \t prima appoggio era solo esterno, mentre nel 63 vicepresidente del consiglio è Nenni (partito socialista) \\
Presidente del consiglio è Aldo moro (d.c.) \\
Diventera nel ?? presidente della repubblica Saragat, che è social democratico \\
Moro rimarra al potere fino al 68  \\
Il centro sinistra pero si affievolisce \t italia ha crisi economica, con calo della produzione e aumento della disoccupazione \\
Nascono contrasti tra partiti di centro (es. d.c.) e partito socialista su come affrontare queste difficolta \\
Vengono fatte varie proposte, che pero sono in contrapposizione \\
Calo della produz. si ha perche c'è sovrapproduzione, dovuta dalla saturazione del mercato \t tutti avevano comprato beni durevoli, che non sono + una novita \\
Conflitti tra partiti, e Ugo La Malfa (ministro del bilancio del governo fanfani) \t ma è del parito repubblicano \\
Centro vuole interenire sull'economia \\
La sinistra non è d'accordo con questa politica, perche la sinistra voleva realizzare servizi sociali + efficienti \t cosi da garantire una sopravvivenza delle classi + deboli \\
Ha la meglio l'idea di un intervento dello stato in campo economico \\
Si aumentano aziende pubbliche, pero meno produttive di quelle private, e diventano degli strumenti clientelari (sono nelle mani del governo)

\v

Serie di riforme millantate non arrivano, e il paese quindi non si ammodernizza veramente \t quindi la parola passa alla contestazione, anche studentesca \t che parte dall'universita di Berkley, in cui gli studenti protestano contro la struttura dell'universita \\
Protesta si allarga all'istruzione, e arriva in europa \t all'italia e poi francia (o controario) \t e in italia passa dai movimenti studenteschi alle fabbriche, e si generano lotte operaie \t ricordate come l'autunno caldo nel 69 \\
Nel 69 proclamato sciopero nazionale del metalmeccanici \t categoria guida della classe operaia \t si vuole rinnovo del contratto di lavoro \\
Poi a novembre sciopero in cui si chiede la casa con scontri, e in dicembre firmato nuovo contratto dei metalmeccanici = vittoria: tutte le richieste vengono accolti
\bi
  \item aumenti salariari per tutti
  \item 40 ore a settimana
  \item diritto ad operai ad organizzare assemblee nelle fabbriche durante orario di lavoro
  \item riduzione delle differenze di trattamento tra operai e impiegati
\ei
Questa richiesta di cambiamento di condizione di lavoro operai, verra accolta anche dal parlamento \t nel 70 approva lo Statuto dei Lavoratori = riconosciuti diritti fondamentali dei lavoratori + fare ricorso al giudice se licenziamento \\
70 è l'anno in cui vengono istituite le regioni, e questi sono gli anni del primo grande referendum: nel 74 c'è referendum sul divorzio, e 60 per cento vota a favore anche se chiesa e d.c. sono contro \\
Poi nell 81 referendum sull aboro \t e solo il 30 per cento sara contrario \t l'influenza sociale della chiesa è ridimensionata \\
Del 68 rimane un cambiamento di mentalita e una direzione verso la laicizzazione dello stato e il valore dell individuo nello stato \t il 68 in quanto rivolta studentesca e movimento operaio si chiude un po li, pero porta cambiamento
\end{document}
